\section{The expansion of \texorpdfstring{$H^2$}{H squared}}\label{exp:section}
The universal bound $H\succeq2(I-\Pi)$ of~\cite{S26} comes from an expansion
of $H^2$ into a triple term and a disjoint-pair term, each bounded
separately. The expansion is due to
Caputo~\cite[Equations~(1.13) to~(1.15)]{Caputo08} and was applied to pair
fibers of graphs in~\cite[Proposition~7.1]{OAI26}. This section computes both terms on the edge indicators of a
regular fiber, where $S$ and $D$ both equal $\tfrac{n-2}2H$, against the
universal bounds $S\succeq0$ and $D\succeq H$, and shows that
on regular graphs the two terms cannot be bounded separately with constants
that add up to the conjectured gap. Nothing here is used in the proofs of
Theorems~\ref{thm:main} and~\ref{thm:small}.

For a vertex triple $T$ let $H_T$ be the sum of $h_a$ over the three pairs
$a\subset T$, and let $D$ be the sum of the products $h_ah_b$ over ordered
pairs $(a,b)$ of disjoint vertex pairs. Put $S=\sum_T(H_T^2-H_T)$. Since
$\sum_TH_T=(n-2)H$, the expansion~\cite[Equation~(6.1)]{S26} reads
\begin{equation}\label{exp:expansion}
 H^2=H+S+D .
\end{equation}
For every graphical degree sequence, $S\succeq0$ and
$D\succeq H$~\cite[Lemma~A.4 and Proposition~6.1]{S26}, which gives
$H^2\succeq2H$. Because $H\succeq0$ has kernel $\mathbb R\mathbf1$, the
equality $\operatorname{gap}(H)=n-1$ on a regular fiber is equivalent to
$H^2\succeq(n-1)H$, that is, to
\begin{equation}\label{exp:target}
 S+D\succeq(n-2)H .
\end{equation}

\begin{proposition}[Edge functions]\label{exp:edges}
Let $\mathcal E$ be the span of $\mathbf1$ and the edge indicators
$\mathbf1[e]$. Each $h_a$ preserves $\mathcal E$ and equals
$\tfrac12(I-T_a)$ there. On $\mathcal E$,
\[
 H=\tfrac12L_T,\qquad H_T^2=\tfrac32H_T\ \text{ for every triple }T,\qquad
 S=\frac{n-2}2H,
\]
and $Dg_{xy}=\tfrac{n-2}2Hg_{xy}=\tfrac{(n-1)(n-2)}2g_{xy}$ for every pair
$\{x,y\}$.
\end{proposition}

\begin{proof}
For $a=\{x,w\}$, the relabeling $T_a$ exchanges $\mathbf1[xy]$ and
$\mathbf1[wy]$ for $y\notin a$ and fixes $\mathbf1[e]$ when $e=a$ or
$e\cap a=\varnothing$. Lemma~\ref{edg:marginal} therefore gives
$E_a=\tfrac12(I+T_a)$ on $\mathcal E$, which proves the first assertion and
$H=\tfrac12L_T$. On $\mathcal E$ we then have
$H_T=\tfrac12(3I-C_T)$, where $C_T$ is the sum of the three transpositions of
the vertices of $T$. Under the group of permutations of $T$, the space $\mathcal E$ is a quotient
of the sum of a trivial representation and the permutation representation on
pairs. The orbits on pairs have size one or three, and
an orbit of size three carries the permutation representation on three
points, which is the sum of the trivial and the two-dimensional irreducible
representation. On these $C_T$ acts as $3$ and as $0$, and the sign
representation does not occur. So $H_T$ has eigenvalues $0$ and $3/2$ on
$\mathcal E$, and $H_T^2=\tfrac32H_T$. Summing $H_T^2-H_T=\tfrac12H_T$ over
$T$ gives $S=\tfrac{n-2}2H$. Finally \eqref{exp:expansion} and
Proposition~\ref{edg:eigen} give
$Dg_{xy}=[(n-1)^2-(n-1)-\tfrac{(n-1)(n-2)}2]g_{xy}$.
\end{proof}

At the functions $g_{xy}$ the two sides of \eqref{exp:target} are equal, and
$S$ and $D$ contribute equally. The next proposition shows that neither
$S\succeq\tfrac{n-2}2H$ nor any other pair of separate bounds summing to
\eqref{exp:target} holds in general.

\begin{proposition}[Obstruction to separate bounds]\label{exp:obstruction}
\begin{enumerate}
\item On $\Omega_{7,2}$ and $T=\{1,2,3\}$, the operator $H_T$ has the
eigenvalue $1$. So the inequality $H_T^2\succeq H_T$
of~\cite[Lemma~A.4]{S26} is attained on a regular fiber, and
$H_T^2\succeq\tfrac32H_T$ fails.
\item On $\Omega_{8,3}$ the number $t$ of triangles of a graph satisfies
\[
 \frac{\langle t,St\rangle}{\langle t,Ht\rangle}=\frac{2314}{835}<3 .
\]
Consequently there are no constants $\alpha,\beta$ with $\alpha+\beta=6$,
$S\succeq\alpha H$ and $D\succeq\beta H$ on $\Omega_{8,3}$.
\end{enumerate}
\end{proposition}

\begin{proof}
For~(1), let $G$ be the $7$-cycle with the cyclic vertex order
$(1,4,2,5,6,3,7)$. No pair inside $T$ is an edge, the vertices $4$ and $7$
each have two neighbors in $T$, and the vertices $5$ and $6$ each have one.
Trades at the pairs inside $T$ preserve these properties and the edge
$\{5,6\}$. The set $\mathcal C$ of graphs reachable from $G$ by such trades
is the set of all assignments of neighborhoods in $T$ to the vertices
$4,5,6,7$ of these sizes in which every vertex of $T$ has degree two; it has
$15$ elements, and functions supported on $\mathcal C$ form an invariant
subspace of $H_T$. Let $c_0$ be the set of graphs in $\mathcal C$ in which
$5$ and $6$ have the same neighbor in $T$, let $c_1$ be the set in which the
neighbor of $5$ in $T$ is adjacent to $4$ and that of $6$ is not, and let
$c_2$ be the set in which the neighbor of $5$ in $T$ is adjacent to $7$ and
that of $6$ is not. These sets have $3$, $6$ and $6$ elements and partition
$\mathcal C$, and they are invariant under the permutations of $T$.
Enumerating the fibers, which have six or two elements, shows that $H_T$
acts on the column vector of coefficients of
$(\mathbf1_{c_0},\mathbf1_{c_1},\mathbf1_{c_2})$ by
\[
 \begin{pmatrix}4/3&-2/3&-2/3\\-1/3&2/3&-1/3\\-1/3&-1/3&2/3\end{pmatrix},
\]
so $f=\mathbf1_{c_1}-\mathbf1_{c_2}$ satisfies $H_Tf=f$. An ancillary program
confirms this identity on $\mathcal C$ in rational arithmetic.
Statement~(2) is a finite computation in rational arithmetic, carried out
by an ancillary program, which evaluates both quadratic forms over the $19{,}355$
elements of $\Omega_{8,3}$; constants do not contribute, since
$H\mathbf1=H_T\mathbf1=0$. The displayed ratio forces $\alpha\le2314/835<3$,
and Proposition~\ref{exp:edges} with $n=8$ forces $\beta\le3$.
\end{proof}
